\section{Finite dictionaries for typical Gaussian networks}
\label{sec:dictionary}
\begin{theorem}[A width-independent finite dictionary]
\label{thm:dictionary}
For $K\ge C(d+H+1)$ and $0<e<1/2$, there is a finite set $\calF_{K,e}$ of functions bounded by $M=C K^{L/2}$ such that
\begin{equation}
 \log|\calF_{K,e}|
 \le C(d+1)e^{-q_H}\log^c\!\bigl(C(d+1)K/e\bigr).
 \label{eq:dictionary-size}
\end{equation}
It depends on $d,H,\phi,K,e$, but not on learner width, prior scale or the measure below. For every probability measure $\xi$ on the input cube, every $N\ge1$ and every scale~\eqref{eq:scales},
\begin{equation}
 \mu_{N,\sigma}\left(\{\|z_w\|_\infty>M\}\cup\left\{
   \inf_{f\in\calF_{K,e}}\|z_w-f\|_{L^2(\xi)}>Me
   \right\}\right)
 \le C e^{-c\omega K}.
 \label{eq:dictionary-prob}
\end{equation}
The exceptional set may depend on $\xi$. The dictionary does not.
\end{theorem}

\subsection{Norm bounds and spectral coordinates}
The event
\begin{equation}
 \frac dN\|W^1\|_F^2\le K,\qquad
 \|W^\ell\|_\op^2\le K\ (2\le\ell\le H),\qquad
 N\|v\|^2\le K
 \label{eq:normevent}
\end{equation}
has probability at least $1-Ce^{-c\omega K}$. For the first and last blocks this is a chi-square bound. For an internal matrix, Gaussian operator-norm concentration around $2\tau$ gives the stated exponent once $K$ exceeds a fixed constant. The first-layer mean energy is $\tau^2(d+1)$, explaining the requirement $K\ge C(d+H+1)$.

On~\eqref{eq:normevent},
\[
 \|h^\ell(x)\|/\sqrt N\le U_\ell,\qquad
 U_1=A\sqrt{K(d+1)/d}+\beta,\qquad
 U_\ell=A\sqrt K\,U_{\ell-1}+\beta.
\]
Hence $U_\ell\le C K^{\ell/2}$ and $|z_w(x)|\le M=C K^{L/2}$.

Fix $\xi$. The normalized hidden kernel
$Q_\ell(x,y)=N^{-1}\langle h^\ell(x),h^\ell(y)\rangle$
is a positive finite-rank operator on $L^2(\xi)$, with trace at most $U_\ell^2$. At a cutoff $\eta_\ell>0$ retain its eigenvalues at least $\eta_\ell$. There is an $N$ by $r_\ell$ matrix $V_\ell$ with orthonormal columns such that
\begin{equation}
 F_\ell(x)=V_\ell^\top h^\ell(x)/\sqrt N,\qquad
 s_\ell(x)=\|(I-V_\ell V_\ell^\top)h^\ell(x)\|/\sqrt N,\qquad
 r_\ell\le U_\ell^2/\eta_\ell.
 \label{eq:state}
\end{equation}
The discarded covariance operator has norm at most $\eta_\ell$.
The pair $(F_\ell,s_\ell)$ has pointwise Euclidean norm at most $U_\ell$.
A Borel choice of spectral bases can be made for the finite Gram matrix, including repeated eigenvalues.

Conditional on earlier layers, the next row preactivation is exactly
\begin{equation}
 a_i^\top F_{\ell-1}(x)+Z_i(x),
 \label{eq:decomposition}
\end{equation}
where the $a_i$ are independent $N(0,\tau^2I)$ and the $Z_i$ are independent centered Gaussian fields, independent of the $a_i$, with covariance $\tau^2$ times the discarded kernel. Their pointwise variance is $\tau^2s_{\ell-1}(x)^2$.
On~\eqref{eq:normevent},
\[
 \frac1N\sum_i a_ia_i^\top\preceq KI,
\]
because the coefficient matrix is $\sqrt N W^\ell V_{\ell-1}$.

\subsection{Removing the Gaussian residual fields}
The next two bounds hold conditionally on all earlier layers and all $a_i$. Write
\[
 h_i(x)=\phi(a_i^\top F(x)+Z_i(x)),\quad
 m_i(x)=\E_Zh_i(x),\quad r_i(x)=h_i(x)-m_i(x).
\]
Suppose the discarded kernel has operator norm at most $\eta$ and pointwise diagonal at most $U^2$.

\begin{lemma}[Residual operator]
\label{lem:residual}
As an operator from $\R^N$ to $L^2(\xi)$, let
$B_rc=N^{-1/2}\sum_i c_i r_i$. Except on conditional probability $e^{-\omega K}$,
\begin{equation}
 \|B_r\|_\op\le C A\left\{\frac U{\sqrt\omega}
                      +\sqrt{\eta(K+1)}\right\}.
 \label{eq:residualop}
\end{equation}
Consequently every orthonormal $N$ by $r$ matrix $V$, even one chosen after the current features are seen, satisfies
\[
 \|V^\top(h-m)/\sqrt N\|_{L^2(\xi;\R^r)}
 \le\sqrt r\,\|B_r\|_\op.
\]
\end{lemma}
\begin{proof}
For fixed $\|c\|=1$, independence and Gaussian Poincar\'e give
$\E\|B_rc\|^2\le A^2\tau^2U^2/N=A^2U^2/\omega$.
As a function of the independent standard Gaussian row variables,
$\|B_rc\|$ is $A\tau\sqrt{\eta/N}$-Lipschitz. Indeed apply pointwise Cauchy--Schwarz in the row index, the Lipschitz property of $\phi$, and the operator norm $\sqrt\eta$ of the residual Gaussian covariance square root into $L^2(\xi)$.

Gaussian concentration~\cite{gaussian}, a $1/2$-net of the unit sphere of cardinality at most $5^N$, and the usual net bound on operator norm prove~\eqref{eq:residualop}. The squared deviation used in the union bound is
$2A^2\eta(K+\tau^2\log5)$: its tail exponent is $\omega K+N\log5$.
The last assertion follows by summing the squared bound over the columns of $V$.
\end{proof}

\begin{lemma}[A uniform residual standard deviation]
\label{lem:rms}
Let $P$ be an orthogonal projector of rank $r$ and put
\[
 T_P(x)=\|(I-P)h(x)\|/\sqrt N,\qquad
 b_P(x)=\sqrt{\|(I-P)m(x)\|^2/N+
                      N^{-1}\sum_i\operatorname{Var}(h_i(x))}.
\]
On the norm event, except on conditional probability at most $Ce^{-\omega K}$, simultaneously for every such projector of every allowed rank,
\begin{align}
 \|T_P-b_P\|_{L^2(\xi)}
 \le C A\biggl\{
 U\sqrt{\frac{r+1}{\omega}}
 +\sqrt{\eta\left[K+(r+1)
                  \log\left(2+\frac{\sqrt{r+1}}{\kappa}\right)\right]}
 +\kappa U_{\rm next}
 \biggr\}
 \label{eq:rms}
\end{align}
for $0<\kappa<1$. Here $U_{\rm next}$ is the next-layer norm bound, and constant factors in it are harmless.
\end{lemma}
\begin{proof}
First fix $P$ independently of the current residual fields. The map from those standard Gaussian variables to $\|T_P-b_P\|_{L^2(\xi)}$ is $A\tau\sqrt{\eta/N}$-Lipschitz. Pointwise Poincar\'e gives
\[
 \operatorname{Var}(T_P(x))
       \le A^2\tau^2s_{\rm prev}(x)^2/N.
\]
The row fluctuations are independent, so
\[
 \E T_P(x)^2=b_P(x)^2-
                 N^{-1}\sum_iP_{ii}\operatorname{Var}(h_i(x)).
\]
The last term is between zero and $rA^2\tau^2s_{\rm prev}(x)^2/N$.
Using $\E(T-\sqrt{\E T^2})^2\le2\operatorname{Var}(T)$ and the square-root inequality gives
\[
 \E\|T_P-b_P\|_{L^2}\le C A U\sqrt{(r+1)/\omega}.
\]

A $\kappa$-net of the rank-$r$ Stiefel manifold in Frobenius norm has cardinality at most $(1+2\sqrt r/\kappa)^{Nr}$. Its projectors approximate in operator norm to within $2\kappa$. On the norm event, $\|h\|/\sqrt N\le U_{\rm next}$; the coefficient Gram bound also gives $\|m\|/\sqrt N\le C U_{\rm next}$. Replacing $P$ by a net projector therefore changes both $T_P$ and $b_P$ in $L^2$ by at most $C\kappa U_{\rm next}$. Gaussian concentration and this net prove~\eqref{eq:rms}. The net contribution to the squared deviation is at most $C\eta\tau^2r\log(2+\sqrt r/\kappa)$, bounded by the displayed expression.

For simultaneous ranks, assign failure probability $e^{-\omega K-r-1}$ at rank $r$ and sum. The additional term in the squared deviation is bounded by $C\eta(r+1)$ and is already covered. Ranks cannot exceed $N$. The norm-event intersection handles the interpolation bounds; its complement was counted separately in~\eqref{eq:normevent}.
\end{proof}

For the actual spectral projector at layer $\ell$, set
$b_i=\sqrt N(V_\ell)_{i,*}$, so $N^{-1}\sum_i b_ib_i^\top=I$.
Let $G_0$ be an independent scalar standard normal and define
\[
 Y_i(x,G_0)=\phi(a_i^\top F_{\ell-1}(x)+\tau s_{\ell-1}(x)G_0).
\]
The two lemmas approximate the actual pair $(F_\ell,s_\ell)$ by
\begin{equation}
 \left(\E_{i,G_0}b_iY_i,\quad
       \sqrt{\E_{i,G_0}Y_i^2-\|\E_{i,G_0}b_iY_i\|^2}\right).
 \label{eq:modelpair}
\end{equation}
Although $b_i$ depends on the current residual fields, the factor bound was uniform over orthonormal matrices and the RMS bound uniform over projectors. Thus this substitution does not assume false conditional independence.

\subsection{Why errors do not acquire repeated square roots}
The coordinate functions of $b_i$ are orthonormal in
$L^2(\text{uniform }i\times\text{Gaussian }G_0)$. For any scalar $Y$, the map
\[
 Y\longmapsto\left(\E bY,\ \|Y-\operatorname{proj}_bY\|_2\right)
\]
is $1$-Lipschitz by Pythagoras and the reverse triangle inequality.
Together with the coefficient Gram bound, changing a previous state $(F,s)$ to $(F',s')$ changes~\eqref{eq:modelpair} by at most
\begin{equation}
 A\sqrt{K\|F-F'\|_{L^2}^2+\tau^2\|s-s'\|_{L^2}^2}.
 \label{eq:propagation}
\end{equation}
No derivative of $\phi$ is used. In particular the argument never differentiates a ReLU Gaussian mean with respect to variance at zero.

Apply Lemma~\ref{lem:sampling} to~\eqref{eq:modelpair}, evaluated on the \emph{already approximated} preceding state, with
\[
 q=1+\|a\|^2/K+\|b\|^2+G_0^2,\qquad
 Z\le2+r_{\rm prev}+r.
\]
Then $Y/\sqrt q=\psi_c(\sqrt K\,a'^\top F+g's)$, where
\[
 \psi_c(t)=c\phi(t/c),\quad c=1/\sqrt q,\quad
 a'=a/(\sqrt K\sqrt q),\quad g'=\tau G_0/\sqrt q .
\]
All normalized parameters have norm at most one; the normalized feature and its second moment are bounded by $C U_\ell$. The vector
$u'=(a',g')$ satisfies the empirical-Gram conditions of Remark~\ref{rem:inputgram}, since its importance-weighted expected Gram is
$\operatorname{diag}(\E aa^\top/K,\tau^2)\preceq I$.
Consequently the selected sampled layer, including Gram whitening and radial projection to radius $U_\ell$, is $C A\sqrt K$-Lipschitz in its incoming state, independently of its ranks.

At the first layer apply the same sampling lemma directly to the empirical distribution of actual incoming rows and $b_i$, with
$q=1+(d+1)\|W^1_{i,*}\|^2/K+\|b_i\|^2$.
There is no Gaussian residual-removal error at that layer.

At the final hidden layer a scalar approximation is enough. Write the readout as $N^{-1}\sum_i c_i h_i$ with $c_i=Nv_i$ and $N^{-1}\sum_i c_i^2\le K$. Lemma~\ref{lem:residual} gives
\[
 \left\|z-\E_{i,G_0}c_iY_i\right\|_{L^2}
 \le C\sqrt K\left\{U_{H-1}/\sqrt\omega+
                         \sqrt{\eta_{H-1}(K+1)}\right\}.
\]
The scalar importance construction after Lemma~\ref{lem:sampling} compresses this expectation. Its selected empirical input Gram and coefficient second moment ensure state Lipschitz constant $C A K$. For $H=1$ apply this scalar construction directly to the first-layer row distribution.

\subsection{Quantitative ranks and a finite dictionary}
We give the bookkeeping because the polynomial degree is part of the claim.
For $H\ge2$ set the final absolute output accuracy to $\varepsilon=C K^{L/2}e$.
Take local state errors
\[
 \delta_\ell=c_H\varepsilon K^{-(H-\ell+1)/2},\qquad
 \eta_{H-1}=c_H\varepsilon^2/K^2,
\]
and choose successively
\[
 \eta_{\ell-1}
 =c_H\frac{\delta_\ell^2\eta_\ell}{K U_\ell^2 J},
 \qquad J=C_H\log(C_HK/e).
\]
Small constants are chosen in depth order to absorb the Lipschitz factors in~\eqref{eq:propagation}. With $\kappa_\ell=c_H\delta_\ell/U_\ell$, the residual errors in Lemmas~\ref{lem:residual} and~\ref{lem:rms} are at most $\delta_\ell$, apart from their finite-effective-width terms.
The deterministic rank caps satisfy
\begin{equation}
 R_\ell\le C_H e^{-2(H-\ell)}J^{H-\ell-1},
                  \qquad 1\le\ell\le H-1.
 \label{eq:rankcaps}
\end{equation}
The importance sample counts can be chosen as
\[
 k_\ell\le C_H(1+R_{\ell-1}+R_\ell)e^{-2}\log^c(C_HK/e),
\]
with $Z_1\le3+R_1$ at the first layer and the corresponding scalar count at the terminal layer. The total scalar parameter count, including the normalized incoming and outgoing coefficient vectors, is bounded by
\begin{equation}
 C_H(d+1)e^{-(4H-2)}\log^c(C_HK/e).
 \label{eq:paramcount}
\end{equation}
For example the largest product is of order $e^{-2}R_1^2$; subsequent layers and the scalar terminal fit the same bound.

The finite-effective-width terms are small whenever
\begin{equation}
 \omega\ge T:=C_H K^{-1}e^{-2}
              \max(1,R_2,\ldots,R_{H-1})\log^c(C_HK/e).
 \label{eq:effectivewidth}
\end{equation}
The first layer contributes no such condition. If $\omega<T$, then $N\le\omega<T$, and we instead cover the original finite network directly on~\eqref{eq:normevent}. The number of parameters is at most
$C_H(T^2+dT)$, which fits~\eqref{eq:paramcount}. Entrywise parameter bounds and network Lipschitz constants are polynomial in $K,T,d$ at fixed depth; parameter quantization adds only logarithms. This small-width argument does not change the prior or introduce a mixture of widths. For $H=1$ the direct scalar sampling argument has no effective-width threshold.

It remains to discretize the scalar perspective functions without assumptions on the asymptotic behavior of $\phi$. Each $\psi_c$ is $A$-Lipschitz and $|\psi_c(0)|\le\beta$. On a bounded interval $[-T_0,T_0]$, the class of all such scalar functions has a uniform $\eta$-net with logarithmic cardinality
\[
 C(1+AT_0/\eta)\log(2+(\beta+AT_0)/\eta).
\]
Construct it by a mesh of spacing $\eta/(CA)$ and quantized vertex values, with piecewise-linear interpolation. The representatives can be taken $2A$-Lipschitz. This factor is absorbed into the depth-dependent constants.

At layer $\ell<H$ the argument interval has radius $T_0\le C U_\ell$.
Choose scalar activation accuracy
$\eta_\ell'=c_H\varepsilon K^{-(H-\ell+1)/2}/\sqrt{Z_\ell}$;
at the terminal layer choose
$\eta_H'=c_H\varepsilon/\sqrt{KZ_H}$.
The empirical input-Gram constraints and the state Lipschitz bounds above propagate these errors to total output error at most $C\varepsilon$. In both cases
\[
 T_0/\eta_\ell'\le C_H\sqrt{Z_\ell}/e.
\]
Thus independently choosing the scalar activation at every atom costs at most
\[
 \sum_\ell C_H k_\ell\sqrt{Z_\ell}/e
 \le C_H e^{-3H}\log^c(C_HK/e)
\]
in log cardinality. For $H\ge2$, $3H\le4H-2$, so~\eqref{eq:paramcount} still dominates.
For $H=1$, bounded $Z$ and $O(e^{-2})$ scalar samples give exponent three.

Finally quantize the remaining bounded program parameters. The ridge is bounded below by a polynomial in $e$ and $K^{-1}$, so matrix inversion, inverse square root and the augmented residual formula have polynomial parameter Lipschitz constants. State radial projections are contractions, and the chosen scalar functions have bounded Lipschitz constants. A sufficiently fine polynomial parameter mesh therefore costs only additional logarithms. Representatives can be selected from nonempty mesh cells intersecting the empirical-Gram constraint set, preserving its state Lipschitz bounds. Clip the terminal output to $[-M,M]$.

The resulting finite set is independent of $\xi,N,\sigma$. Conditional failure probabilities from the Gaussian lemmas sum over the fixed number of layers to $Ce^{-c\omega K}$. Intersecting with~\eqref{eq:normevent}, choosing small depth-dependent approximation constants, and the direct small-width cover prove Theorem~\ref{thm:dictionary}.
